\section{Follow-the-Sun：全球值班不等于自动拥有清晰责任}

长训练跨越时区，follow-the-sun 可以减少无人响应时间，但只有 handoff packet、明确 owner 和 stop criteria 才能避免重复操作。Frontier training 的难点是系统和研究责任交织：网络故障由 infra owner 处理，loss drift 需要 research owner，data corruption 又跨 pipeline 与 privacy owner。

\subsection{Handoff packet 与 escalation}

一个有效 handoff 至少包含当前 checkpoint/step、expected curve、最近异常、已验证 hypothesis、禁止重复的 action、pending experiment、下一决策 deadline 与 contacts。接班者必须 acknowledge 并复核 dashboard，而不是只读聊天记录。High-risk action，例如修改 optimizer 或跳过数据 shard，应要求双人 review。

\lecturefigure{08-follow-the-sun.png}{Follow-the-sun 依赖 active owner、handoff packet、接班确认与跨职能 escalation。}{本地字幕 19:00--21:20；概念重绘。}

\begin{knowledgebox}{读图：Coverage 与 ownership 是两件事}
全球 coverage 保证有人在线，ownership 决定谁有权 stop、rollback 或 change recipe。若多个地区都能“顺手修”，run 会积累不可追踪变更；若所有人只观察，又会错过 recovery window。Handoff 应冻结当前 hypothesis 和 authority，重大变更写入 run ledger。
\end{knowledgebox}

\teachervoice{讲者承认 follow-the-sun 也没有消除困难：有些 failure 模糊、偶发且跨研究/系统边界。课堂提醒我们，增加时区覆盖不是组织魔法，仍需共享语言、runbook 与明确 escalation。}

\subsection{本章小结}

全球训练运营的关键是 evidence transfer 与 decision rights。Handoff 质量决定 coverage 能否转化为更低 downtime，而不是更多并发误操作。

